\usepackage{amsmath}                                    % 排版数学公式
\usepackage{latexsym}
\usepackage{amsfonts}                                   %数学符号字体库宏包套件，它包含有：amsfonts、amssymb、eufrak 和 eucal 四个宏包。
\usepackage{amssymb}                                    % 定义AMS的数学符号命令
\usepackage{mathrsfs}                                   % 数学RSFS书写字体
\usepackage{bm}                                         % 数学黑体
\usepackage{graphicx}                                   % 支持插图,图形宏包graphics的扩展宏包
\usepackage{float}                                      % 提供 [H] 浮动体精确定位（图片用），原先由 minted 间接引入

\usepackage{color}                                      % 支持彩色
\usepackage[dvipsnames]{xcolor}
\usepackage{amscd}
\usepackage[linesnumbered,ruled,vlined]{algorithm2e}
\usepackage{diagbox}
\usepackage{titlesec}                                   %设置章节格式
\usepackage{enumerate}                                 	%更改enumerate环境格式
\usepackage{hyperref}

\usepackage{ccicons}

\newtheorem{axiom}{公理}
\newtheorem{theorem}{定理}
\newtheorem{corollary}{推论}

\hypersetup{
    colorlinks=true,
    linkcolor=Blue,
}

\raggedbottom

% --- Simplified-Chinese-only footnote --------------------------------------
% \scfootnote{...} typesets a footnote ONLY in the Simplified Chinese (Mainland)
% edition; in the uncensored original and in every translation it renders nothing
% (the argument is silently dropped). The boolean below is OFF here; the only thing
% that turns it on is the Simplified Chinese build, through a single censor.csv rule
% that rewrites the toggle line -- censor.csv being, by design, the sole difference
% between /text and that edition. See scripts/languages/Simplified Chinese/prompt.md.
\newif\ifSCversion
\SCversionfalse
\newcommand{\scfootnote}[1]{\ifSCversion\footnote{#1}\fi}
